\slajdid{02-s003}
\begin{frame}[label=02-s003]{Ulazni signal u logičko kolo se menja u vremenu}
\gusttekst
Signal – fizička veličina koja nosi informaciju o logičkim jedinicama i nulama.\\
Za sada još uvek razmišljamo na apstraktnom nivou sa logičkim jedincama i nulama

\centering
\begin{tikzpicture}[DEoznaka,baseline=(current bounding box.center)]
\node[not gate US,draw,minimum width=6mm,minimum height=5mm] (g) at (0,0) {};
\draw[DEvod] (g.input)--++(-.5,0) node[left] {\(x(t)\)};
\draw[DEvod] (g.output)--++(.5,0) node[right] {\(y(t)\)};
\end{tikzpicture}\par
\begin{columns}[c,onlytextwidth]
\begin{column}{.40\textwidth}\centering
\begin{tikzpicture}[DEoznaka,x=10mm,y=7.2mm]
\draw[DEstrelica] (-.7,1.9)--(3.8,1.9) node[below] {\(t\)};
\draw[DEstrelica] (0,1.75)--(0,3.7) node[left] {\(x(t)\)};
\draw[DEstrelica] (-.7,-.1)--(3.8,-.1) node[below] {\(t\)};
\draw[DEstrelica] (0,-.25)--(0,1.7) node[left] {\(y(t)\)};
\node[above left,inner sep=1pt] at (0,3) {=1};\node[above left,inner sep=1pt] at (0,2.1) {=0};
\node[above left,inner sep=1pt] at (0,1) {=1};\node[above left,inner sep=1pt] at (0,.1) {=0};
\draw[DEvod] (-.5,2.1)--(.6,2.1)--(.6,3)--(1.6,3)--(1.6,2.1)--(3.5,2.1);
\draw[DEvod] (-.5,1)--(.6,1)--(.6,.1)--(1.6,.1)--(1.6,1)--(3.5,1);
\draw[dashed] (0,3)--(.6,3);
\draw[dashed] (.6,2.1)--(.6,-.25) node[below] {\(t_1\)};
\draw[dashed] (1.6,2.1)--(1.6,-.25) node[below] {\(t_2\)};
\draw[DEcrvena] (-.8,-.65)--(3,3.5) (-.8,3.5)--(3.5,-.65);
\end{tikzpicture}

\end{column}
\begin{column}{.42\textwidth}\centering
\begin{tikzpicture}[DEoznaka,x=10mm,y=7.2mm]
\draw[DEstrelica] (-.7,1.9)--(3.8,1.9) node[below] {\(t\)};
\draw[DEstrelica] (0,1.75)--(0,3.7) node[left] {\(x(t)\)};
\draw[DEstrelica] (-.7,-.1)--(3.8,-.1) node[below] {\(t\)};
\draw[DEstrelica] (0,-.25)--(0,1.7) node[left] {\(y(t)\)};
\node[above left,inner sep=1pt] at (0,3) {=1};\node[above left,inner sep=1pt] at (0,2.1) {=0};
\node[above left,inner sep=1pt] at (0,1) {=1};\node[above left,inner sep=1pt] at (0,.1) {=0};
\draw[DEvod] (-.5,2.1)--(.6,2.1)--(.6,3)--(1.6,3)--(1.6,2.1)--(3.5,2.1);
\draw[DEvod] (-.5,1)--(1.15,1)--(1.15,.1)--(2.5,.1)--(2.5,1)--(3.5,1);
\draw[dashed] (0,3)--(.6,3) (0,.1)--(1.15,.1);
\draw[dashed] (.6,2.1)--(.6,-.3) (1.6,2.1)--(1.6,-.3);
\draw[dashed] (1.15,.1)--(1.15,-.3) (2.5,.1)--(2.5,-.3);
\draw[<->] (.6,-.25)--(1.15,-.25) node[midway,below] {\(t_{pHL}\)};
\draw[<->] (1.6,-.25)--(2.5,-.25) node[midway,below] {\(t_{pLH}\)};
\draw[DEcrvena,<->] (.6,2.3)--(1.6,2.3);
\draw[DEcrvena,<->] (1.15,.85)--(2.2,.85);
\node[DEcrvena] at (1,1.45) {Uočiti};
\draw[DEcrvena,->] (1.05,2.2)--(1.5,.98);
\end{tikzpicture}

\end{column}
\begin{column}{.16\textwidth}\gusttekst
p – propagation\\
H – high\\
L – low
\end{column}
\end{columns}
\raggedright
\(t_{pHL}\) - kašnjenje izlaznog signala u odnosu na promenu ulaznog signala prilikom prelaska IZLAZNOG signala sa „nivoa“ H na L

\(t_{pLH}\) - kašnjenje izlaznog signala u odnosu na promenu ulaznog signala prilikom prelaska IZLAZNOG signala sa „nivoa“ L na H\\
H = 1, L = 0
\end{frame}
